\section{Introduction}\label{sec:introduction}

Polynomial identity testing asks whether a polynomial given by arithmetic
operations is zero. In the noncommutative setting, variables are evaluated
on matrices: scalar substitutions would, for example, erase the nonzero
commutator $x_1x_2-x_2x_1$. The question considered here is whether one can
construct a single matrix substitution, using only the characteristic, the
number of variables, and a formula-size bound, that detects every nonzero
formula of that size.
We construct such a substitution over the prime field in every positive
characteristic. Its dimension and the time needed to print all its entries
are polynomially bounded.

\subsection{The result}

For a field $F$, the free associative algebra
$F\langle x_1,\ldots,x_n\rangle$ has as an $F$-basis the finite ordered
words in $x_1,\ldots,x_n$, including the empty word $1$; multiplication
concatenates words. A \emph{noncommutative formula} is a rooted tree whose
leaves are variables or scalars in $F$, and whose internal vertices are
binary addition or ordered multiplication gates. Its \emph{size} is the
number of all vertices, including scalar leaves. Thus there are no inverse
gates and no shared subexpressions. Scalars may be arbitrary elements of $F$.
When evaluating at $d\times d$ matrices, a scalar $a$ is interpreted as $aI_d$.

\begin{theorem}\label{thm:main}
There is a deterministic Turing machine and absolute positive integers $C,k$ with
the following property. Given a prime $p$ in binary and positive integers
$n,s$ in unary, the machine outputs a positive integer $d=O(n^3s^6)$ and
a tuple
\[
 (T_1,\ldots,T_n)\in\Mat_d(\F_p)^n.
\]
For every field $F$ of characteristic $p$ and every nonzero
$f\in F\langle x_1,\ldots,x_n\rangle$ computed by a formula of size at
most $s$, the same tuple, viewed over $F$, satisfies
\[
 f(T_1,\ldots,T_n)\ne0.
\]
Writing $L=\lceil\log_2p\rceil$, the total running time and output length
are at most $C(n+s+L+1)^k$. All matrix entries are output densely as ordinary
binary residues in $\{0,\ldots,p-1\}$, with dimensions and separators included.
Primality of $p$ is promised.
\end{theorem}

The generator receives no formula or coefficients. In particular, its
guarantee is uniform over arbitrary extensions of $\F_p$, and applies to
$p=2$ without any relation between the characteristic and the formula size.
The required conclusion is matrix nonvanishing. The time to evaluate a
formula over an unspecified coefficient field is separate from the bit
complexity of constructing the tuple.

The proof also applies to an \emph{acyclic algebraic path program}: a finite
directed acyclic graph with a designated source and sink, edges labelled by
scalars or variables, and value equal to the sum of ordered label products
over source--sink paths. Replacing the formula-size parameter by the number
of vertices gives the same polynomial construction; see
Corollary~\ref{cor:path-programs}. This extension follows from the coefficient
representation used in the proof.

\subsection{Context and the characteristic obstruction}

Noncommutative identity testing has distinct tests that use the internal
representation and tests that only evaluate the polynomial. Nisan
characterized homogeneous noncommutative branching-program complexity by
ranks of coefficient matrices~\cite[Theorem~1]{Nisan1991}. Raz and Shpilka gave
a deterministic polynomial-time identity test for noncommutative formulas
and algebraic branching programs when their descriptions are
available~\cite{RS2005}. Matrix evaluation also underlies randomized
black-box tests; Bogdanov and Wee gave such a test for circuits with a
supplied degree bound, using polynomial identities of matrix
algebras~\cite[Theorem~4.2]{BW2005}.

For deterministic black-box testing, Forbes and Shpilka constructed
quasipolynomial-size matrix hitting sets for noncommutative algebraic
branching programs, under a field-cardinality hypothesis
\cite[Corollary~4.7]{FS2013}. A block-diagonal sum can combine a finite
hitting set into one tuple, but its dimension grows with the size of that
set. Lakhani and Mukhopadhyay constructed one integer matrix tuple for
sparsity-bounded polynomials over characteristic-zero fields
\cite[Theorem~2.1 and Remark~2.6]{LM2026}; their finite-field formulation
also distinguishes a succinct extension-field description from its explicit
coordinate cost. Sparsity is a different parameter: a formula of small size
can have exponentially many words. These results provide relevant
benchmarks, while the construction here is polynomial in formula size
and in the binary length of the characteristic.

The characteristic-zero companion~\cite[Theorem~1.1]{OAInc} constructs one
rational matrix tuple for formulas, with dimension at most $2ns^2$ and
polynomial bit complexity. Its proof encodes words by iterated integrals,
puts the encoding of a formula into a finite differential system, and
applies an order-of-vanishing bound before truncating. The representation
by an acyclic path program is independent of the characteristic. The
analytic-looking part of the construction is not: formal integration
divides by positive integers, ordinary differentiation annihilates
$p$th powers in characteristic $p$, and distinct integer pole locations
may coincide. Nor does reduction of a rational hitting tuple automatically
preserve its guarantee at every prime.

We retain the word-series and matrix-system organization, replacing the
differential operator by an infinite-order multiplicative shift. The
resulting order bound and the final descent to prime-field matrices are
proved directly in the given characteristic. The dimension bound is larger
than in the companion, but remains polynomial without any lower bound on $p$.

\subsection{Proof strategy}

Fix an arbitrary coefficient field $F$ and introduce an indeterminate $t$
over it. Over $K=F(t)$, let $\sigma$ act on formal series in $z$ by
$\sigma(g)(z)=g(tz)$. For each letter $i$, define the operator $\J_i$ by
requiring $\J_i(g)$ to have zero constant term and
\[
 (\sigma-1)\J_i(g)=\frac{z}{1-(1+t)^i z}\,g.
\]
On the coefficient of $z^r$ for $r\ge1$, this divides by $t^r-1$, which is nonzero
in every characteristic. Iteration defines a series $h_w$ for each word
$w$, starting from $h_\eps=1$ for the empty word $\eps$.
Section~\ref{sec:words} proves that these series are linearly independent.
A minimal-relation argument reduces the claim to a pole obstruction:
a difference $\sigma(R)-R$, with $R\in K(z)$, cannot have just one pole
in an infinite shift orbit,
whereas the poles $(1+t)^{-i}$ lie in disjoint such orbits.
Thus a nonzero polynomial $f=\sum_w f_wx_w$ gives a
nonzero series $h_f=\sum_w f_wh_w$.

Formula size bounds how late the first nonzero coefficient of $h_f$ can
appear. Section~\ref{sec:systems} finds a row $u$, a column $v$, and
strictly upper triangular $m\times m$ matrices $B_1,\ldots,B_n$ over $F$,
with $m\le2s$, such that $f_w=uB_wv$; here
$B_{i_1\cdots i_\ell}=B_{i_1}\cdots B_{i_\ell}$ and $B_\eps=I_m$.
The finite sum $H=\sum_w B_wh_w$ satisfies $H(0)=I$ and
$\sigma(H)=(P/Q)H$, with $z$-degrees at most $n$ and $Q(0)=1$.
Section~\ref{sec:order} proves the reusable bound
\[
 uHv\ne0\quad\Longrightarrow\quad
 \ord_z(uHv)\le n\binom m2.
\]
The proof adapts the determinant method for differential systems: a basis
of series with distinct leading orders $c_i$ has a shifted determinant
whose leading coefficient is a Vandermonde determinant on $t^{c_i}$.
These elements remain distinct in small characteristic. Rational degree
bounds then control the orders of the relevant minors.

Truncating in $z$ now gives matrices over $\F_p(t)$ that detect all the
required formulas. Section~\ref{sec:construction} removes $t$ without
choosing a large finite extension. Strict triangularity supplies one
denominator clearing every evaluated word, together with a polynomial
bound on the remaining $t$-degree. Reduction modulo a sufficiently large
power $t^E$ preserves nonvanishing. Finally, representing multiplication
on $\F_p[t]/(t^E)$ by matrices over $\F_p$ gives the output tuple.
The shift order bound is useful independently of the formula application;
the last step is a general way to retain bounded-degree rational data in
ordinary matrices over the base field.
